%==============================================================================
\part{Translation}
%==============================================================================

\section{The arithmetic of partial identifier translation}
\label{sec:translation}

Every federated plan that crosses a database boundary crosses a namespace
boundary, and the crossing is performed by a cross-reference map. These maps
are the least examined component of federated retrieval and the one where the
losses actually accumulate. This section treats them as first-class objects of
the language and derives what a plan can and cannot know about a chain of them.

\subsection{Maps are partial, non-injective, and non-confluent}

\begin{definition}[Translation map]
\label{def:map}
A \emph{translation map} from namespace $n$ to namespace $n'$ is a relation
$\mu \subseteq n \times n'$. We write $\mu(S) = \{ v : u \in S,\ (u,v) \in \mu\}$
for the image of a set, $\operatorname{dom}\mu = \{u : \exists v,\ (u,v) \in \mu\}$,
and call $\mu$ \emph{total} if $\operatorname{dom}\mu = n$, \emph{functional}
if each $u$ has at most one image, and \emph{injective} if each $v$ has at most
one preimage.
\end{definition}

Real cross-reference maps are none of the three. A chemical entity may have no
counterpart in a pathway database (partiality); a single ontology class may
correspond to several database compounds differing only in protonation state
(non-functionality); and several distinct ontology classes may map to one
database compound that does not resolve the distinction (non-injectivity).

\begin{definition}[Retention]
\label{def:retention}
The \emph{retention} of $\mu$ on a set $S$ is
\[
r_\mu(S) \;=\; \frac{|S \cap \operatorname{dom}\mu|}{|S|}
\qquad (S \neq \emptyset),
\]
the fraction of $S$ that survives translation, and $r_\mu(\emptyset) = 1$ by
convention. The \emph{amplification} is
$a_\mu(S) = |\mu(S)| / |S \cap \operatorname{dom}\mu|$ when the denominator is
nonzero.
\end{definition}

Retention and amplification are independent: a map may retain everything and
amplify (each input has several images), retain little and not amplify, or
both. Reporting only the output cardinality $|\mu(S)|$ conflates them, which is
why \cref{def:retention} separates them and why the prototype records both.

\begin{proposition}[Cardinality alone is uninformative]
\label{prop:cardinality-uninformative}
For any $N \geq 1$ and any output size $M \geq 1$ there exist maps $\mu_1,\mu_2$
and a set $S$ with $|S| = N$ such that
$|\mu_1(S)| = |\mu_2(S)| = M$ while $r_{\mu_1}(S) = 1$ and
$r_{\mu_2}(S) = 1/N$.
\end{proposition}

\begin{proof}
Let $S = \{u_1,\dots,u_N\}$. For $\mu_1$, distribute $M$ images across all $N$
inputs so that every input has at least one when $M \geq N$, and when $M < N$
let the inputs share images so that all $N$ lie in the domain and the image has
$M$ elements; then $r_{\mu_1}(S)=1$. For $\mu_2$, let
$\operatorname{dom}\mu_2 \cap S = \{u_1\}$ and let $\mu_2(u_1)$ have $M$
elements; then $|\mu_2(S)| = M$ and $r_{\mu_2}(S) = 1/N$. The output
cardinalities agree and the retentions differ by a factor of $N$.
\end{proof}

\Cref{prop:cardinality-uninformative} is the reason a plan cannot detect a
translation failure by looking at how much came out. In \cref{lst:script} the
translation is line~7 and its output size is the only thing the script observes
about it; by the proposition, that observation carries no information about
whether the translation lost anything. This is defect (D3), and it is not
fixable by checking that the output is non-empty.

\subsection{Composition}

\begin{theorem}[Retention of a chain]
\label{thm:retention}
Let $\mu_1 : n_0 \to n_1, \dots, \mu_k : n_{k-1} \to n_k$ and let $S_0$ be
finite and non-empty, with $S_i = \mu_i(S_{i-1})$. Then
\begin{enumerate}[label=(\alph*),leftmargin=2.4em]
\item the end-to-end retention factorises multiplicatively along the realised
sets,
\[
\frac{|S_k|}{|S_0|}
\;=\;
\prod_{i=1}^{k} r_{\mu_i}(S_{i-1})\, a_{\mu_i}(S_{i-1}),
\]
with the convention that the product is $0$ if any $S_{i-1}$ is empty;
\item if each $\mu_i$ is injective on the realised sets, the \emph{surviving
fraction} --- the fraction of $S_0$ having at least one image in $S_k$ ---
satisfies
\[
\rho_{1..k}(S_0) \;\leq\; \min_{1 \leq i \leq k} r_{\mu_i}(S_{i-1}),
\]
and the bound is attained when each $\mu_i$ is in addition functional. Without
injectivity the inequality fails, and $\rho_{1..k}(S_0) \leq r_{\mu_1}(S_0)$ is
all that survives (\cref{rem:injectivity-needed});
\item under the same injectivity hypothesis,
$\rho_{1..k}(S_0) \geq 1 - \sum_{i=1}^{k}\bigl(1 - r_{\mu_i}(S_{i-1})\bigr)$,
and the bound is attained when the loss sets are pairwise disjoint under the
realised correspondence. Injectivity is again load-bearing: without it a single
lost element of $S_{i-1}$ can extinguish several survivors of $S_0$ at once, and
the bound fails (\cref{rem:injectivity-needed}).
\end{enumerate}
\end{theorem}

\begin{proof}
(a) By \cref{def:retention}, $|S_i| = |\mu_i(S_{i-1})| =
a_{\mu_i}(S_{i-1}) \cdot |S_{i-1} \cap \operatorname{dom}\mu_i| =
a_{\mu_i}(S_{i-1}) r_{\mu_i}(S_{i-1}) |S_{i-1}|$ whenever $S_{i-1}$ is
non-empty. Telescoping over $i = 1..k$ and dividing by $|S_0|$ gives the
product. If some $S_{i-1}$ is empty then $S_k$ is empty and both sides are
zero.

(b) An element of $S_0$ survives to $S_k$ only if it lies in
$\operatorname{dom}\mu_1$, and its image lies in $\operatorname{dom}\mu_2$, and
so on. Let $A \subseteq S_0$ be the surviving set and fix a stage $i$. Each
$u \in A$ has at least one trajectory reaching a surviving element of
$S_{i-1}$, so the map sending $u$ to one such element is a function
$A \to S_{i-1} \cap \operatorname{dom}\mu_i$. If every $\mu_j$ is injective on
the realised sets this function is itself injective, whence
$|A| \leq r_{\mu_i}(S_{i-1})|S_{i-1}| \leq r_{\mu_i}(S_{i-1})|S_0|$ and the
minimum bound follows. Adding functionality gives each element exactly one
trajectory, and taking all losses to occur at the single stage achieving the
minimum attains the bound. The case $i=1$ needs no hypothesis, since survival
requires membership of $\operatorname{dom}\mu_1$ and distinct elements of $S_0$
are distinct there; this is the residual bound
$\rho_{1..k}(S_0) \leq r_{\mu_1}(S_0)$.

(c) The set of elements of $S_0$ that fail to survive is contained in the union
over $i$ of the sets $D_i \subseteq S_0$ whose trajectory dies at stage $i$.
A trajectory dies at stage $i$ by reaching an element of $S_{i-1}$ outside
$\operatorname{dom}\mu_i$, and there are $(1-r_{\mu_i}(S_{i-1}))|S_{i-1}|$ such
elements. Under injectivity on the realised sets, distinct members of $D_i$
reach distinct elements of $S_{i-1}$ --- the same injection constructed in (b)
--- so $|D_i| \leq (1-r_{\mu_i}(S_{i-1}))|S_{i-1}| \leq
(1-r_{\mu_i}(S_{i-1}))|S_0|$, the last step because injectivity also forces
$|S_{i-1}| \leq |S_0|$. A union bound over the $k$ stages gives the stated
lower bound, with equality when the $D_i$ are pairwise disjoint. Without
injectivity $|D_i|$ is not controlled by $|S_{i-1}|$ at all, since one dead
element of $S_{i-1}$ may be the sole continuation of arbitrarily many members
of $S_0$.
\end{proof}

\begin{remark}[Injectivity is a hypothesis of (b) and (c), not a nicety]
\label{rem:injectivity-needed}
The hypothesis in \cref{thm:retention}(b) and (c) cannot be dropped. Take
$S_0 = \{u_1,\dots,u_9\}$ with $\mu_1$ retaining seven of them, sending $u_9$
to $\{t_9\}$ and $u_{10}$ to $\{t_9, t_{10}\}$ --- non-injective, but with
amplification $a_{\mu_1}(S_0) = 1$, so the collision is invisible in the
factorisation of (a). Let $\mu_2$ retain three of the seven, among them $t_9$.
Then $r_{\mu_1} = 7/9$, $r_{\mu_2} = 3/7$, and four elements of $S_0$ survive,
giving $\rho = 4/9 > 3/7 = \min_i r_{\mu_i}$. Two survivors in $S_0$ are riding
one survivor in $S_1$, so the surviving set upstream is larger than the
surviving set downstream and the minimum is not an upper bound for it.

The failure is non-injectivity specifically, not amplification: $a_{\mu_1} = 1$
throughout the example. A cross-reference table that maps two source entries
onto one target entry --- routine wherever one resource merges what another
distinguishes --- is exactly this configuration, so the hypothesis is a real
restriction on which chains the two bounds may be applied to, and a plan
reporting coverage should measure $\rho$ rather than bound it.

The lower bound (c) fails for the same reason, by a different route. Its union
bound counts elements of $S_0$ whose trajectory dies at stage $i$ by counting
the elements of $S_{i-1}$ that lie outside $\operatorname{dom}\mu_i$; when
several survivors of $S_0$ have been merged onto one element of $S_{i-1}$, a
single such loss extinguishes all of them, and the count understates the
damage. A two-stage instance: $|S_0| = 8$ with realised sizes
$|S_1| = 6$, $|S_2| = 5$ and stage retentions $3/4$ and $5/6$, so (c) predicts
$\rho \geq 1 - 1/4 - 1/6 = 7/12$, while the measured surviving fraction is
$1/2$. The one element lost at stage two carried two survivors, costing $2/8$
where the bound had budgeted $1/6$. Note that the na\"ive repair ---
requiring $|S_{i-1}| \leq |S_0|$ --- does not help, since it already holds here.

Sampling two-stage chains over eight elements with each element retained
independently with probability $0.8$ and given up to three images, and keeping
only the non-injective draws, (b) fails for $33.9\%$ of chains and (c) for
$2.0\%$ ($n = 5000$). Under injective draws neither fails, as the theorem now
requires.
\end{remark}

\begin{remark}[The two bounds are far apart, and the gap is the point]
\label{rem:bounds-gap}
For $k=3$ and $r_{\mu_i} = 0.8$ throughout, on an injective chain, (b) gives
$\rho \leq 0.8$ and (c) gives $\rho \geq 0.4$. The true value depends on
whether the three maps lose \emph{the same} entities or \emph{different} ones,
and no amount of pairwise coverage data determines which. A plan that needs to
know its own coverage must therefore measure the chain, not compute it from
published per-map figures --- which is what \cref{def:retention-check} makes it
do, and what \cref{prop:pairwise-insufficient} says it cannot avoid.
\end{remark}

\begin{proposition}[Pairwise coverage does not determine chain coverage]
\label{prop:pairwise-insufficient}
There exist two families $(\mu_1,\mu_2)$ and $(\mu_1',\mu_2')$ with
$r_{\mu_i} = r_{\mu_i'}$ for $i=1,2$ on a common $S_0$, whose end-to-end
surviving fractions differ.
\end{proposition}

\begin{proof}
Let $S_0 = \{a,b,c,d\}$, and let both first maps be the identity on
$\{a,b,c\}$ and undefined on $d$, so $r_{\mu_1} = r_{\mu_1'} = 3/4$. Let
$\mu_2$ be the identity on $\{a,b\}$ and undefined on $c$, and let $\mu_2'$ be
the identity on $\{b,c\}$ and undefined on $a$; both have retention $2/3$ on
$S_1 = \{a,b,c\}$. The chain through $\mu_2$ retains $\{a,b\}$, a surviving
fraction of $2/4$; the chain through $\mu_2'$ retains $\{b,c\}$, also $2/4$. To
separate them, extend $\mu_2'$ to be defined on $a$ as well and undefined on
both $b$ and $c$: then $r_{\mu_2'} = 1/3 \neq 2/3$, so instead take
$S_0 = \{a,b,c,d,e,f\}$, let $\mu_1 = \mu_1'$ be the identity on
$\{a,b,c,d\}$ (retention $4/6$), let $\mu_2$ be the identity on $\{a,b\}$ and
$\mu_2'$ the identity on $\{c,d\}$, each of retention $2/4$. Both chains
survive $2$ of $6$. Now instead let $\mu_1'$ be the identity on
$\{c,d,e,f\}$ --- still retention $4/6$ --- with $\mu_2'$ the identity on
$\{c,d\}$ of retention $2/4$; the primed chain still survives $\{c,d\}$. Finally
take $\mu_2''$ to be the identity on $\{e,f\}$, retention $2/4$ on
$S_1' = \{c,d,e,f\}$: the chain $\mu_1'$ then $\mu_2''$ survives $\{e,f\}$,
whereas the chain $\mu_1$ then $\mu_2''$ has $S_1 = \{a,b,c,d\}$ meeting
$\operatorname{dom}\mu_2'' = \{e,f\}$ in the empty set, surviving fraction $0$.
The two chains have identical pairwise retentions $(4/6, 2/4)$ and end-to-end
surviving fractions $2/6$ and $0$.
\end{proof}

\subsection{Order does not recover loss}

A natural hope is that a plan losing too much to translation can be repaired by
reordering: translate later, filter earlier, cross the boundary at a different
point. The hope is false in a precise sense.

\begin{theorem}[Reordering does not recover translation loss]
\label{thm:reorder}
Let $\Plan$ contain a translation step $\mu$ and a source step $\sigma$ whose
denotation is a selection $\pi$, and let $\Plan'$ be the plan obtained by
exchanging their order where both orders are well formed. Then
\[
\mu(\pi(S)) \subseteq \mu(S) \cap \pi'(\mu(S))
\quad\text{and}\quad
\pi'(\mu(S)) \subseteq \mu(S),
\]
where $\pi'$ is the image selection. In particular neither order retains an
element of $S$ that lies outside $\operatorname{dom}\mu$, so the two orders
have the same upper bound $r_\mu(S)$ on surviving fraction, and reordering
cannot increase it.
\end{theorem}

\begin{proof}
Both containments are immediate from monotonicity of images and selections.
For the substantive claim, let $u \in S \setminus \operatorname{dom}\mu$. In
the order ``select then translate'', $u$ either fails $\pi$ and is dropped, or
passes $\pi$ and then has no image under $\mu$; either way it contributes
nothing to the output. In the order ``translate then select'', $u$ has no image
under $\mu$ and so contributes nothing before $\pi'$ is applied. Hence in both
orders the output is contained in $\mu(S \cap \operatorname{dom}\mu)$, whose
preimage in $S$ has size at most $r_\mu(S)|S|$. Since this bound is independent
of the order, no exchange increases it.
\end{proof}

\begin{remark}[What reordering does change]
\label{rem:reorder-cost}
\Cref{thm:reorder} says reordering cannot recover \emph{coverage}. It says
nothing about \emph{cost}, which it changes substantially: translating after a
restrictive selection sends fewer identifiers across the boundary and so costs
fewer requests, whereas translating first may allow a batched request where the
other order does not. The planner is therefore free to reorder for cost, and
\cref{thm:allocation} governs that freedom --- but a plan that reorders hoping
to retain more is mistaken, and the language does not offer the operation as a
coverage remedy.
\end{remark}

\subsection{Routes may disagree without either being wrong}

\begin{definition}[Route]
\label{def:route}
A \emph{route} from namespace $n$ to namespace $n'$ is a path in the directed
graph whose vertices are namespaces and whose edges are available translation
maps. Two routes are \emph{parallel} if they share their endpoints.
\end{definition}

\begin{theorem}[Route divergence measures resolved extent]
\label{thm:route-extent}
Let $R = \mu_k \circ \cdots \circ \mu_1$ and $R' = \mu'_l \circ \cdots \circ \mu'_1$
be parallel routes from $n$ to $n'$, each composed of maps that are
\emph{sound} --- every pair they assert is a true correspondence --- and let
$S \subseteq n$. Then:
\begin{enumerate}[label=(\alph*),leftmargin=2.4em]
\item $R(S)$ and $R'(S)$ cannot contradict: every element of each is a true
correspondent of some element of $S$, so neither set contains an element the
other must exclude;
\item the symmetric difference $R(S) \bigtriangleup R'(S)$ consists exactly of
correspondences resolved by one route and unresolved by the other, so
$|R(S) \bigtriangleup R'(S)|$ is a lower bound on the number of correspondences
that at least one route fails to resolve;
\item if every map on both routes is total on the realised sets then
$R(S) = R'(S)$, so divergence implies partiality somewhere;
\item consequently $R(S) \cup R'(S)$ is sound and has surviving fraction at
least $\max(\rho_R(S), \rho_{R'}(S))$.
\end{enumerate}
\end{theorem}

\begin{proof}
(a) By soundness of each constituent map and transitivity of true
correspondence, every $v \in R(S)$ stands in true correspondence to some
$u \in S$, and likewise for $R'$. A set of true correspondents cannot
contradict another such set: contradiction would require one route to assert a
pair the other denies, and neither route denies anything --- both are
existential.

(b) Let $v \in R(S) \setminus R'(S)$. By (a), $v$ is a true correspondent of
some $u \in S$, and $R'$ does not produce it, so $R'$ fails to resolve the
correspondence $(u,v)$. Symmetrically for the other difference. Each element of
the symmetric difference therefore witnesses at least one unresolved
correspondence on one route, and distinct elements witness distinct
correspondences, giving the bound.

(c) If every map is total on the realised sets, each route computes the full
composite correspondence relation restricted to $S$, and the composite relation
is determined by the underlying true correspondence, which is route-independent.
So the two images coincide.

(d) The union of two sound sets is sound, and contains each, so its surviving
fraction is at least each of theirs.
\end{proof}

\begin{remark}[The union is not free]
Part (d) suggests always taking the union of parallel routes. The cost is that
the union costs the sum of the routes' requests, and by
\cref{thm:allocation} that budget comes from somewhere. The language exposes
\texttt{union} over routes precisely so that the trade is written down rather
than assumed, and \cref{ex:route-union} shows the syntax.
\end{remark}

\begin{example}[Two routes, written down]
\label{ex:route-union}
\begin{lstlisting}[caption={Parallel routes made explicit. The plan asserts nothing about which is better; it measures the divergence and reports it.},label=lst:routes]
let direct   = map compounds via chebi_to_kegg
let indirect = map compounds via chebi_to_inchikey
               then via inchikey_to_kegg

let both     = union direct indirect
assert soundness(direct) and soundness(indirect)
emit divergence(direct, indirect) as resolved_extent
\end{lstlisting}
By \cref{thm:route-extent}(b) the emitted \texttt{resolved\_extent} is a lower
bound on the correspondences at least one route misses. It is not an error
count and must not be reported as one: both routes may be individually correct
and the divergence is then a statement about coverage, not about accuracy.
\end{example}

\subsection{The retention check}

\begin{definition}[Retention check]
\label{def:retention-check}
A translation step annotated \texttt{expect partial} $\epsilon$ computes
$r_\mu(S)$ on its actual input and, if $r_\mu(S) < \epsilon$, returns
$\vstarve$ to its successors with a diagnosis naming $\mu$, $\epsilon$ and the
observed $r_\mu(S)$. Otherwise it returns $\vans$ with the image and records
both $r_\mu(S)$ and $a_\mu(S)$ as provenance attributes.
\end{definition}

\begin{proposition}[The check converts a silent decimation into a located verdict]
\label{prop:check-locates}
Let $\Plan$ contain a translation step $\Step_t$ with retention below its
declared $\epsilon$, feeding a step $\Step_d$ that would otherwise report
$\vempty$. Under \cref{def:retention-check}, $\Step_d$ reports $\vstarve$ and
the blame chain of \cref{prop:blame} terminates at $\Step_t$.
\end{proposition}

\begin{proof}
By \cref{def:retention-check}, $\Step_t$ returns a non-$\vans$ verdict. By (R1)
of \cref{def:verdicts}, $\Step_d$ reports $\vstarve$ naming $\Step_t$. By
\cref{prop:blame} the chain terminates at a non-$\vstarve$ step; since
$\Step_t$'s verdict was assigned by \cref{def:retention-check} rather than by
(R1), and $\Step_t$ names no predecessor, the chain terminates there.
\end{proof}

This is the resolution of defect (D3) in its commonest concrete form. Without
the check, the analyst sees an empty pathway result and investigates the
pathway database. With it, the executor reports that the translation two steps
earlier retained $0.31$ against a declared $0.6$, and names the map.
